\documentclass[10pt,a4paper]{amsart}
\usepackage[margin=19mm]{geometry}
\usepackage{amsmath,amssymb,mathtools}
\usepackage[scr=rsfs,cal=euler]{mathalfa}
\usepackage{xfrac}
\usepackage[unicode,hidelinks,bookmarks=false]{hyperref}
\newcommand{\ie}{i.e.\ }
\newcommand{\cf}{cf.\ }
\newcommand{\resp}{respectively\ }
\newcommand{\Subsection}{\subsection}
\providecommand{\nobreakdash}{\nobreak\hbox{-}}
\setlength{\emergencystretch}{2em}
\multlinegap0pt
\usepackage{amssymb}
\usepackage{tikz}
\usepackage{enumerate}
\usepackage{multirow}
\usepackage{booktabs}
\usepackage{xcolor}
\newtheorem{thm}{Theorem}
\newcommand{\Sr}{S_{\vec{r}}}
\newcommand{\Z}{\mathbb{Z}}
\title{Grouped Stirling complexes}
\author{Alessia Revelli, Steven Scheirer}
\date{Source: arXiv:2602.09884v1 (CC BY 4.0)}
\begin{document}
\maketitle

\noindent\emph{Source and licence.} Grouped Stirling complexes. Version arXiv:2602.09884v1 (CC BY 4.0), distributed by arXiv under the Creative Commons Attribution 4.0 International licence, http://creativecommons.org/licenses/by/4.0/, as recorded in the arXiv licence metadata for this exact version and shown on the abstract page https://arxiv.org/abs/2602.09884v1. Full licence notice: \texttt{licenses/arxiv-2602.09884-CC-BY-4.0.txt}.\par\smallskip

\noindent\textit{Editorial scope.} Counted unit: the article's thm thm (source0.tex lines 3508--3517). The article's complete original proof of it is delivered with it (source0.tex lines 3519--3548, one \texttt{proof} environment of 30 lines). No mathematics is written, repaired or re-derived: the statement and the proof are the article's own lines, in the article's own notation, and no retained line is substituted or re-flowed.  Retained uncounted as the source-local context the proof needs: lines 3489--3491.  Nothing was omitted from any retained span, so the package carries no omission record.  No retained line contains a citation, so the wrapper carries no bibliography and no bibliography entry.  No retained line contains a figure environment, an image-inclusion command or drawing code, so no image is delivered and the package declares no asset dependency.  Editorial formatting only, in addition: an AMS \texttt{amsart} preamble that restates the article's own declarations and macros used by the retained lines, namely the environments \texttt{thm} on separate counters, together with the standard packages those lines need.  The retained environments print in this document as Theorem 1 in order of appearance, whereas the article prints the same items with the numbering of its own full document.  The complete native source of the article as distributed in the arXiv e-print for this exact version is bundled unchanged as \texttt{sources/194351/source0.tex}.
\begin{equation}\label{eqn:total}
    \binom{r}{i} m^in^{r-i}.
\end{equation}

\begin{thm}\label[thm]{thm:numCellsB}
    Let $G$ be a simple graph with $n$ vertices and $m$ edges, and let $\vec{r}=(n-1,n-1,\dots,n-1)\in\Z^r$ for some integer $r\ge 2$. In $\Sr(G)$ we have 
    \[
    \#(\text{i-cells)}=\begin{cases}\displaystyle
        \binom{r}{i}\left(m^in^{r-i}-\displaystyle\sum_{v\in V(G)}\deg(v)^i\right)&\text{if }0\le i<r\\
        m^r+m-\displaystyle\sum_{v\in V(G)}\deg(v)^r&\text{if }i=r\\
        0&\text{if }i> r\\
    \end{cases}.
    \]
\end{thm}

\begin{proof}
As noted above, to count the number of $i$-cells of $\Sr(G),$ it suffices to count the number of valid $i$-edge-$r$-tuples by subtracting the number of invalid $i$-edge-$r$-tuples from (\ref{eqn:total}).

If $\tilde c$ is an invalid $i$-edge-$r$-tuple, then any vertex in $\tilde c$ must be a $k$-missed vertex for \textit{every} vertex-only color $k$. This means that each vertex in $\tilde c$ is equal to some fixed vertex $v.$ Furthermore, every edge in $\tilde c$ must have $v$ as an endpoint, since if $k$ is an edge-color and $e$ is the $k$-colored edge, and $v$ is not an endpoint of $e$, there would be  a $k$-colored robot on $v$ (and hence $\tilde c$ would be valid).


Therefore, every invalid $i$-edge-$r$-tuple is of the form $\tilde c=(x_1,x_2,\dots,x_r),$ where there is some fixed vertex $v$ where $x_k=v$ when $x_k$ is a vertex, and $x_k$ has $v$ as an endpoint if $x_k$ is an edge. Thus for each $v\in V(G),$ we get an invalid $i$-edge-$r$-tuple $\tilde c$ by choosing $i$ of the $r$ components for the edges, then for each of these components, choosing one of the $\deg(v)$ edges which have $v$ as an endpoint. The remaining components of $\tilde c$ must be equal to $v.$ So, for each fixed vertex $v,$ we get $\binom{r}{i}\deg(v)^i$ invalid $i$-edge-$r$-tuples in which $v$ appears.

Furthermore, for $i<r,$ any $i$-edge-$r$-tuple must contain at least one vertex, so that in this case, distinct vertices give distinct invalid $i$-edge-$r$-tuples, so, the total number of invalid $i$-edge-$r$-tuples is 
\begin{equation}\label{eqn:invalid1}
    \sum_{v\in V(G)}\binom{r}{i}\deg(v)^i.
\end{equation}
Subtracting (\ref{eqn:invalid1}) from (\ref{eqn:total}) gives the result for $i<r.$

For $i=r$, any invalid $r$-edge-$r$-tuple $\tilde c$ must be of the form $\tilde{c}=(x_1,\dots,x_r),$ where each $x_k$ is an edge, with a common endpoint $v.$ So for each vertex $v,$ we must count the number of $r$-tuples of edges where each edge has $v$ as an endpoint. As in the previous case, there are $\deg(v)$ choices for each of the $r$ edges, so there are a total of $\deg(v)^r$ invalid $r$-edge-$r$-tuples in which all edges have $v$ as an endpoint. Consider the sum 
\begin{equation}\label{eqn:double}
    \sum_{v\in V(G)}\deg(v)^r.
\end{equation}

Since an invalid $r$-edge-$r$-tuple consists exclusively of edges, for any edge $e$ with endpoints $v_1$ and $v_2$, the invalid $r$-edge-$r$-tuple $\tilde c = (e,e,\dots,e)$ is double-counted in (\ref{eqn:double}): once in the term corresponding to $v_1$ and once in the term corresponding to $v_2.$  If $\tilde c$ consists of at least two different edges, then the edges in $\tilde c$ only have one common endpoint, so $\tilde c$ is only counted once. So, the only invalid $r$-edge-$r$-tuples which are double-counted are of the form $(e,e,\dots,e).$ Since there are $m$ edges in $G$, there are $m$ $r$-edge-$r$-tuples which are double-counted.

Therefore, the total number of invalid $r$-edge-$r$-tuples is 
\begin{equation}\label{eqn:Invalidrr}
    \left(\sum_{v\in V(G)}\deg(v)^r\right)-m.
\end{equation}

Subtracting (\ref{eqn:Invalidrr}) from (\ref{eqn:total}) (with $i=r$) gives the result.

The fact that there are no $i$-cells for $i>r$ follows from the fact that for each color $k$, there is at most one $k$-colored edge in any given cell $c$, so there are at most $r$ edges in total in $c$.
\end{proof}

\end{document}
